% SSU Data Dash 2026: team report template. Compile from this folder (Overleaf: upload the whole
% repository as a zip). `python run.py report` writes the results table and figure it includes.
% Keep the report to 4 pages, not counting the outside-data table and the appendix.
\documentclass[11pt]{article}
\usepackage[margin=2.2cm]{geometry}
\usepackage{amsmath,graphicx,booktabs,hyperref}
\graphicspath{{../figures/}}
\setlength{\parskip}{4pt}

% Required: your team name and the full name of every team member.
\title{SSU Data Dash 2026: \textit{Duo Sucker}}
\author{\textit{Preston Cai, Lorance Jiang}}
\date{}

\begin{document}
\maketitle

\noindent\begin{tabular}{@{}ll}
\textbf{Video (YouTube):} & \url{https://youtu.be/YOUR-VIDEO} \\
\textbf{Code and submission files (GitHub):} & \url{https://github.com/Preston-Cai/SSU-Data-Dash} \\
\end{tabular}

\section*{Summary}
We built separate ensembles for off-market and list-aware valuation.  Each combines
boosted trees using property and strictly historical local-sale features with a
TF--IDF text model that reads listing remarks.  In six-month walk-forward
validation (December 2025--May 2026), the current rerun produced the results in
Table~1.  Our main finding is that carefully timed local-sale history matters
more than changing tree libraries: a neighbourhood feature is useful only when
it is constructed from sales already known on the prediction date.

\section{Data and features}
% Which columns you used, what you built from them, and anything you dropped and why.
We used physical attributes (area midpoints, lot size, bedrooms, bathrooms,
parking, style, and property type), municipality and community, rounded
latitude/longitude cells, calendar and seasonal variables, and indicators
extracted from remarks.  Local-history features summarize log price and log
price per square foot for community/type/style and geographic groups.  The
list-aware track additionally uses list price and historical sold-to-list
ratios; the off-market pipeline never receives list price.  We excluded the
external AVM and fields known only after sale.  No outside data were used.

We also supplied the trees with log area and log lot area, area-per-bedroom
and bathrooms-per-bedroom ratios, garage/parking indicators, condominium
flags, multi-resolution geographic cells, and community--property-type and
community--style interactions.  Remarks became compact renovation, condition,
luxury, transit, basement, and ``as-is'' indicators; the separate Ridge model
kept the full word-level remarks rather than reducing them to flags.

The targets were chosen to match the metric: off-market models predict
$\log(\text{sold\_price})$, while list-aware models predict
$\log(\text{sold\_price}/\text{list\_price})$.  This reduces the influence of
the largest homes and makes the list-aware task a percentage adjustment rather
than a second full-price estimate.  Local statistics are smoothed toward the
overall market when a group is rare.  We retained both ordinary encodings and
recency-weighted encodings because the former estimates a stable local level,
while the latter can respond to market drift.  The final text training window
is limited to the latest 36 months, and the strongest list-aware recency
encoder uses a four-month half-life.

\section{Models}
% One paragraph per track. Off-market has no list price; list-aware does. Explain how your two
% models differ and why.
The off-market prediction is a calibrated geometric blend of recency-aware
LightGBM (32\%), word TF--IDF Ridge (19\%), and XGBoost (49\%).  The list-aware
blend uses ordinary target-encoded LightGBM (11.3\%), recency-aware LightGBM
(37.7\%), an independent off-market value anchor (15.5\%), and a smoothed
historical community-plus-property-type sold-to-list median (35.5\%).  Trees capture nonlinear interactions
among structure, location, and earlier comparable sales; Ridge contributes a
different view through listing language.  We blend them because their errors
are not identical, even when a component is weaker by itself.  The specialized
list-aware anchor helps decide how much to trust an asking price by comparing it
with both prior local sale ratios for the same property family and a valuation
that never saw list price.

We also tested a dense PyTorch MLP with standardized numeric features,
property/style indicators, and learned municipality/community embeddings.  It
reached 11.19\% off-market and 5.63\% list-aware MdAPE, so we did not promote it.
Its feature set lacked remarks, richer interactions, and the historical local
price signals that most helped the tree ensemble; this particular architecture
also fit mixed tabular structure less effectively than boosted trees.

The screening experiments explain the final model choice.  On the November
screen, leakage-safe local encodings reduced off-market MdAPE from 8.86\% to
7.75\%, recency weighting reached about 7.69\%, and XGBoost reached 7.62\%.
The components were retained because their errors were complementary.  In
contrast, CatBoost reached 10.34\% off-market MdAPE, a direct comparable
estimator reached 11.83\%, and a local-anchor residual model reached about
7.85\%.  These alternatives were therefore rejected rather than included for
complexity alone.  A weighted geometric blend was selected because the models
operate on log-price targets; a multiplicative calibration then removes a
remaining systematic bias.

\section{Validation}
% How you tested the models without using information from the future: which months you
% trained on, which you predicted, and how you made sure nothing leaked.
We used walk-forward validation rather than a random split.  For each month from
December 2025 through May 2026, every component was rebuilt using only sales
strictly before that month's cutoff, then evaluated on that month.  Target
encodings are also chronological \emph{within} the training set: a training
home's local statistics and global smoothing prior use only sales with an
earlier \texttt{date\_sold}.  Homes sold on the same date are processed as a
batch and cannot reveal their targets to one another.  Validation homes use the
complete permitted pre-cutoff history, and their sold prices never enter feature
construction.  Blend weights and calibration were fitted on
September--November 2025.  The reported list-aware result is the conservative
3.98\% score from the pre-specified walk-forward evaluation; later
post-validation tuning is not used as evidence for the final result.
Across the six official months, monthly MdAPE ranged from 6.98\% to 7.63\%
off-market and from 3.64\% to 4.08\% list-aware; May was the weakest
off-market month, consistent with market drift rather than a single anomalous
row.

\section{Results}
\begin{table}[h]
\centering
\input{../reports/results_table.tex}
\caption{Walk-forward validation on the listings. MdAPE: median absolute percentage error.
MAPE: mean absolute percentage error. Brackets: bootstrap 95\% interval.}
\end{table}
The table is regenerated directly from row-level predictions by
\texttt{python run.py report}.  MdAPE describes the typical percentage miss,
while MAPE exposes larger errors.  We report both because a model can improve
the median while leaving an unacceptable tail.  The rerun is the authoritative
score: earlier experimental numbers used leave-one-out encodings that did not
meet the handout's stricter ``earlier sales only'' rule and are not claimed here.

The saved final artifact contains the fitted encoders, component models, blend
weights, calibration, and submission files.  Reloading that checkpoint
reproduces the predictions byte-for-byte, and both CSVs pass the starter
submission checker.  This separates the reproducible submission model from
the exploratory screening scripts.

\section{Where the model fails}
\begin{figure}[h]
\centering
\includegraphics[width=0.85\linewidth]{error_by_price_range.png}
\caption{Absolute percentage error by sold-price decile.}
\end{figure}
Errors rise sharply in the expensive tail: the highest sold-price decile is
about 13\% MdAPE off-market versus about 6--7\% in the cheapest deciles, while
list-aware error rises to about 5\% from roughly 3\%.  Sparse communities and
rare property combinations also fail more often because their historical
encodings are pulled toward the market prior.  The off-market task is harder
because it lacks the seller's asking price; the month-by-month range above also
shows sensitivity to market drift.  Text helps when remarks reveal renovation
or condition, but synthetic or vague descriptions add little.  These patterns
suggest calibrated uncertainty by price and local-history support, not just a
single point estimate.

\section{Limitations}
The June--September hidden test remains the only truly unseen evaluation, so
walk-forward improvement is evidence rather than a guarantee.  Rare groups are
still noisy despite smoothing, and recency weighting cannot anticipate a sudden
market shift.  The ensemble is also slower and harder to audit than a single
model.  December--May was inspected during development, so future tuning should
use nested time splits or a newly reserved period.  Next we would add
segment-level uncertainty intervals and test whether simpler components can be
removed without sacrificing accuracy.

\section*{Outside data}
% Every outside source, or write "None". Data may only cover periods up to and including May 2026.
\begin{tabular}{lll}
\toprule
Source & URL & Periods covered \\
\midrule
\textit{None} & & \\
\bottomrule
\end{tabular}

\section*{Reproducing our results}
% The exact commands that produce your two submission files from a fresh copy of your repository.
\texttt{pip install -r gbt\_ensemble\_experiments/requirements-attempts.txt}\\
\texttt{python gbt\_ensemble\_experiments/validate.py}\\
\texttt{python run.py report}\\
\texttt{python gbt\_ensemble\_experiments/train\_final.py}\\
\texttt{python run.py check}

\end{document}
